\section{Local Search}

\textit{Local search} bekerja pada satu (atau sekumpulan) \textit{state} lengkap dan
berpindah ke \textit{state} tetangga untuk memperbaiki nilai fungsi objektif. Jalur
pencarian tidak disimpan, sehingga kebutuhan memorinya kecil dan cocok untuk
persoalan optimasi berukuran besar seperti persoalan pengisian peti kemas ini.
Kelemahannya, pencarian dapat terjebak di \textit{local optimum}, \textit{plateau},
atau \textit{ridge}~\cite{russell2010aima,khodra2021localsearch}.

\section{Hill-Climbing}

Hill-Climbing selalu berpindah ke tetangga yang lebih baik dan berhenti ketika tidak
ada tetangga yang lebih baik~\cite{russell2010aima,khodra2021hc}. Varian yang
diimplementasikan:
\begin{itemize}
	\item \textbf{Steepest-ascent}: membangkitkan semua tetangga dan memilih yang terbaik.
	Berhenti bila tetangga terbaik tidak lebih baik daripada \textit{state} sekarang.
	\item \textbf{Stochastic}: memilih secara acak salah satu tetangga yang lebih baik.
	\item \textbf{Sideways move}: seperti \textit{steepest-ascent}, tetapi boleh berpindah ke
	tetangga yang bernilai sama untuk keluar dari \textit{plateau}, dibatasi oleh
	parameter \textit{maximum sideways move}.
	\item \textbf{Random-restart}: menjalankan \textit{steepest-ascent} berulang kali dari
	\textit{state} awal acak yang berbeda dan mengambil hasil terbaik, dibatasi oleh
	parameter \textit{maximum restart}.
\end{itemize}

\section{Simulated Annealing}

Simulated Annealing~\cite{kirkpatrick1983sa,khodra2021sa} memilih satu tetangga acak
setiap iterasi. Tetangga yang lebih baik ($\Delta E>0$) selalu diterima, sedangkan
tetangga yang lebih buruk diterima dengan peluang $e^{\Delta E/T}$. Suhu $T$
diturunkan bertahap sesuai jadwal (\textit{schedule}). Pada suhu tinggi pencarian
bersifat eksploratif, sedangkan pada suhu rendah perilakunya mendekati Hill-Climbing.
Dengan demikian SA dapat keluar dari \textit{local optimum} di awal pencarian.

\section{Genetic Algorithm}

Genetic Algorithm~\cite{russell2010aima,khodra2021ga,kkif2026slides} memelihara
populasi berisi $k$ \textit{state} (individu). Nilai setiap individu diukur dengan
\textit{fitness function}. Generasi berikutnya dibentuk melalui tiga operator:
\begin{enumerate}
	\item \textbf{Seleksi}: induk dipilih dengan peluang sebanding \textit{fitness}
	(\textit{roulette wheel}).
	\item \textbf{Crossover}: dua induk dipotong pada titik acak yang sama, lalu
	bagian ekornya ditukar sehingga terbentuk dua anak.
	\item \textbf{Mutasi}: anak diubah secara acak dengan peluang kecil.
\end{enumerate}
